
In order to provide preliminar validation of the proposed design, we simulated a simplified model of our robot with V-REP simulation framework~\cite{vrep13}. Here, we kept the size of vehicle: all motors are displaced on a 160mm radius circle, 60 degrees apart from each other, and tilted 55 degrees away from their vertical axis (see Table~\ref{tab:selected}), but with different mass and inertia, as we do not know yet have a detail design of the mass distribution. In V-REP we simulated the force and torque produced by each propeller as modeled in~\eqref{eq:Fi} and~\eqref{eq:Mi}. Thus, this simulation aims at the validation of both the design discussed in Section~\ref{sec:design-parameters} and the controller proposed in Section~\ref{sec:posit-attit-contr}.

This validation was divided into two steps: on the first, each motion mode (translation and attitude) is simulated, first, the position controller, having it move 1 meter on each direction, and then the attitude controller, by moving 70 degrees on each of its axis; on the second, we simulated a path defined by a sequence of waypoints with attitude setpoints, thus combining both controllers.

\subsection{Individual motion mode}

% The analysis of each mode enables a better adjustment of the gains of the controller and validates each controller, and component, individually. 
For this validation, the robot is holding on position $\mathbf{P_h} = (0,0,1)$, while the reference is set at $\mathbf{P_r} = (1,1,2)$. The total movement of the robot is 1 meter on each direction. On Figure~\ref{fig:position_errors} we provide the obtained results for position convergence with the proposed controller in~\eqref{eq:fctrl} and~\eqref{eq:pctrl}.

\begin{figure}
  \centering
  \subfloat[]{\includegraphics[width=0.33\linewidth]{px-error.png}}
  \subfloat[]{\includegraphics[width=0.33\linewidth]{py-error.png}}
  \subfloat[]{\includegraphics[width=0.33\linewidth]{pz-error.png}}
  \caption{Position Controller Error Convergence: (a)~plot of error on X component, (b)~plot of error on Y component, and (c)~plot of error on Z component.}
  \label{fig:position_errors}
\end{figure}

The three error plots are approximately equal due to the chosen gains, which are equal on all axis for this experiment: the force applied on each component is only dependent on the mass of the vehicle and position error, and not o the moment of inertia.

For the attitude controller different gains are set according to the moment of inertia of the vehicle, which resulted in a faster response for the X axis then for the Z or Y axis. The attitude of the vehicle starts as $\mathbf{R_h} = (0,0,0)$ and the reference is set as $\mathbf{R_r} = {(70\degree,0,0), (0,70\degree,0), (0,0,70\degree)}$, for rotations along X, Y and Z, respectively. The results are shown in Figure~\ref{fig:attitude_errors}.

\begin{figure}
  \centering
  \subfloat[]{\includegraphics[width=0.33\linewidth]{rx-error.png}}
  \subfloat[]{\includegraphics[width=0.33\linewidth]{ry-error.png}}
  \subfloat[]{\includegraphics[width=0.33\linewidth]{rz-error.png}}
  \caption{Attitude Controller Error Convergence: (a)~plot of error on X component, (b)~plot of error on Y component, and (c)~plot of error on Z component.}
  \label{fig:attitude_errors}
\end{figure}

\subsection{Waypoint navigation}

To validate the joint position and attitude control, we used waypoint navigation: we created a virtual path composed by 6 setpoints, varying on both position and attitude. The trajectory followed by the vehicle is represented on \ref{fig:trajectory}. Plotted are both the X and Y axes of the body frame~$\mathcal{B}$ to represent its attitude. We retrieved the errors across the whole trajectory following, for both position and velocity, as well as attitude and angular velocity. These are shown in Figure~\ref{fig:nav_errors}. A video showing this simulation is available online here:~\texttt{https://youtu.be/BW1Q3lryb9c}


\begin{figure}
  \centering
  \includegraphics[width=0.8\linewidth]{3d-trajectory-names.png}
  \caption{Trajectory followed by the robot. Plotted are X and Y axis of the inertia frame~$\mathcal{I}$.}
  \label{fig:trajectory}
\end{figure}


\begin{figure}
  \centering
  \subfloat[]{\includegraphics[width=0.45\linewidth]{traj-p-error.png}} \hfill
  \subfloat[]{\includegraphics[width=0.45\linewidth]{traj-v-error.png}} \\
  \subfloat[]{\includegraphics[width=0.45\linewidth]{traj-r-error.png}} \hfill
  \subfloat[]{\includegraphics[width=0.45\linewidth]{traj-om-error.png}}
  \caption{Waypoint navigation position and attitude errors: (a)~error on position, (b)~error on linear velocity, (a)~error on rotation, and (b)~error on angular velocity.}
  \label{fig:nav_errors}
\end{figure}

%%% Local Variables:
%%% mode: latex
%%% TeX-master: "main"
%%% End:
